\section{Data and Sample}\label{sec:data}

The paper uses the universe of BEC item-transactions from January~2016 to December~2019: 3.7~million bid-level observations across 82{,}569 distinct items, 832{,}984 purchase orders, and 1{,}344 public buyer units. SEFAZ/SP (the state finance secretariat) manages the platform and archives the bidding records. BEC covers all standardized-goods procurement in the state, so the data are free of the selection concerns that accompany voluntary platforms.

Two samples organize the empirical analysis. The reduced-form difference-in-differences in Section~\ref{sec:rf} uses the full BEC universe: Group 65 as the treated cohort, the 76 product groups balanced in pre and post periods as never-treated controls, and an 18-month window around the March~2018 cutoff (Stata monthly-date units $[680, 715]$). This window is symmetric; nested windows of 6 and 12 months deliver the same qualitative findings. The structural estimation in Sections~\ref{sec:model}--\ref{sec:results} narrows the sample to Group-65 Preg\~ao auctions inside that same 18-month window. The restriction is motivated: Preg\~ao is the descending-clock format that admits point identification of cost distributions from drop-out bids, and Group~65 is the only product group that switched regime during the sample. The structural sample is 297{,}967 firm-auction observations across 97{,}993 distinct auctions, balanced across Pre (48{,}740 auctions) and Post (49{,}253) periods.

Bids within the structural sample are subject to three filters. First, the normalized bid $c_\epsilon = b / p^{\mathrm{ref}}$ must satisfy $c_\epsilon \in (0, 3]$; the upper cutoff removes implausible outliers (less than 0.4 percent of bids, almost always data-entry errors at the hundred-times-reference level). Second, auctions must have at least two participating firms. Third, each bid carries a type flag (SME or non-SME, from the BEC administrative SME registry, validated against historical bidding to avoid misclassifying firms that temporarily exited SME status) and a pharma flag (CADMAT classes 6531, 6532, 6536, 6581---CMED-regulated pharmaceuticals and biologicals---are pharma-narrow; the remainder of Group~65 is non-pharma). Descriptive counts appear in Appendix~\ref{app:sample}.

Auction-level unobserved heterogeneity is material in the structural sample and not fully absorbed by the reference-price normalization. Two items in the same CADMAT class can differ in scale, ANVISA registration burden, or volume-specific cost elasticities in ways the reference price does not fully reflect. I handle this with a method-of-moments variance decomposition followed by best-linear-predictor shrinkage \`a la \citet{krasnokutskaya2011}. The shrinkage recovers intraclass correlations $\rho^k$ of 0.47 (non-pharma non-SMEs) to 0.62 (pharma non-SMEs), in the 0.5--0.7 range that Krasnokutskaya (2011) reports for California highway contracts. The UH-clean bids enter all downstream structural estimation; statistical details are in Section~\ref{sec:identification}.

Entry counts are constructed directly. For each auction, I count the distinct firms of each type that submitted at least one bid, then average within (period $\times$ pharma) cells. In Preg\~ao non-pharma auctions, participation moves from 0.94 SMEs and 2.68 non-SMEs per auction pre-policy to 1.87 SMEs and 1.50 non-SMEs post-policy. Pharma moves from 0.55 and 2.61 to 1.22 and 1.66. The doubling of SME participation and the one-third to one-half decline in non-SME participation are the two stylized facts the extensive-margin story has to explain.

The furosemide vignette from the introduction anchors these aggregates in a single item. Item 110639, furosemide 40~mg coated tablets, is a CMED-regulated generic in class 6531. In the pre-period open-regime purchase of February~2018, ten firms competed. In the post-period restricted purchase of October~2018, three SMEs competed, and the winning price cleared just above the reference. The pair illustrates both the extensive and intensive margins: fewer firms, and the three that showed up priced differently than the ten had. Disentangling which of these two margins drives the aggregate price movement is the task of Sections~\ref{sec:rf} onward.

Four ancillary inputs anchor the analysis. Reference prices come from the BEC catalog (\texttt{preco\_ref}) and are used both for normalization and for translating structural estimates into reais. CADMAT class codes come from the BEC product hierarchy. The historical SME flag is reconstructed from firm-level bidding histories over the full 2014--2019 period to protect against transient status changes. Winner-level characteristics (firm size, age, two-digit CNAE sector) come from linking BEC winners to Brazil's matched employer-employee registry (RAIS) for 2017; the link rate is 82.6 percent. All of these are described in detail in Appendix~\ref{app:sample}.
